\documentclass[tikz,border=8pt]{standalone}
% Shared look for every TikZ diagram. Load with % Shared look for every TikZ diagram. Load with % Shared look for every TikZ diagram. Load with \input{../../tools/tikz-style.tex}
\usepackage[T1]{fontenc}
\usepackage{lmodern}
\renewcommand{\sfdefault}{lmr}  % diagrams use the same serif as the text
\usetikzlibrary{arrows.meta, positioning, shapes.geometric, calc, fit, backgrounds}
\definecolor{cblue}{HTML}{4C78A8}
\definecolor{corange}{HTML}{F58518}
\definecolor{cgreen}{HTML}{54A24B}
\definecolor{cred}{HTML}{E45756}
\definecolor{cpurple}{HTML}{B279A2}
\definecolor{cgrey}{HTML}{6B6B6B}
\tikzset{
  every picture/.style={font=\sffamily, line width=1.2pt},
  box/.style={draw=#1, fill=#1!12, rounded corners=4pt, align=center, inner sep=7pt, minimum height=1.1cm},
  box/.default=cgrey,
  arrow/.style={-{Stealth[length=3mm]}, draw=cgrey, line width=1.4pt},
  title/.style={font=\sffamily\bfseries\Large},
  note/.style={font=\sffamily\small, text=cgrey, align=center},
}
\AtBeginDocument{\hyphenpenalty=10000 \exhyphenpenalty=10000}

\usepackage[T1]{fontenc}
\usepackage{lmodern}
\renewcommand{\sfdefault}{lmr}  % diagrams use the same serif as the text
\usetikzlibrary{arrows.meta, positioning, shapes.geometric, calc, fit, backgrounds}
\definecolor{cblue}{HTML}{4C78A8}
\definecolor{corange}{HTML}{F58518}
\definecolor{cgreen}{HTML}{54A24B}
\definecolor{cred}{HTML}{E45756}
\definecolor{cpurple}{HTML}{B279A2}
\definecolor{cgrey}{HTML}{6B6B6B}
\tikzset{
  every picture/.style={font=\sffamily, line width=1.2pt},
  box/.style={draw=#1, fill=#1!12, rounded corners=4pt, align=center, inner sep=7pt, minimum height=1.1cm},
  box/.default=cgrey,
  arrow/.style={-{Stealth[length=3mm]}, draw=cgrey, line width=1.4pt},
  title/.style={font=\sffamily\bfseries\Large},
  note/.style={font=\sffamily\small, text=cgrey, align=center},
}
\AtBeginDocument{\hyphenpenalty=10000 \exhyphenpenalty=10000}

\usepackage[T1]{fontenc}
\usepackage{lmodern}
\renewcommand{\sfdefault}{lmr}  % diagrams use the same serif as the text
\usetikzlibrary{arrows.meta, positioning, shapes.geometric, calc, fit, backgrounds}
\definecolor{cblue}{HTML}{4C78A8}
\definecolor{corange}{HTML}{F58518}
\definecolor{cgreen}{HTML}{54A24B}
\definecolor{cred}{HTML}{E45756}
\definecolor{cpurple}{HTML}{B279A2}
\definecolor{cgrey}{HTML}{6B6B6B}
\tikzset{
  every picture/.style={font=\sffamily, line width=1.2pt},
  box/.style={draw=#1, fill=#1!12, rounded corners=4pt, align=center, inner sep=7pt, minimum height=1.1cm},
  box/.default=cgrey,
  arrow/.style={-{Stealth[length=3mm]}, draw=cgrey, line width=1.4pt},
  title/.style={font=\sffamily\bfseries\Large},
  note/.style={font=\sffamily\small, text=cgrey, align=center},
}
\AtBeginDocument{\hyphenpenalty=10000 \exhyphenpenalty=10000}

% The tree of derivatives of f = x^3 + xy + y^2. Differentiating twice gives four second partial derivatives; the two
% mixed paths (x then y, y then x) both give 1. The four results fill the Hessian.
% Idea after Khan Academy, "Symmetry of second partial derivatives" and "The Hessian matrix"; our own function.
\begin{document}
\begin{tikzpicture}[every node/.append style={font=\large}]
  \node[box=cgrey] (f) at (0,0) {$f = x^3 + xy + y^2$};
  \node[box=cblue] (fx) at (-3.4,-2.3) {$\dfrac{\partial f}{\partial x} = 3x^2 + y$};
  \node[box=cgreen] (fy) at (3.4,-2.3) {$\dfrac{\partial f}{\partial y} = x + 2y$};
  \node[box=cblue] (fxx) at (-5.6,-5) {$\dfrac{\partial^2 f}{\partial x^2} = 6x$};
  \node[box=corange] (fxy) at (-1.75,-5) {$\dfrac{\partial^2 f}{\partial y\,\partial x} = 1$};
  \node[box=corange] (fyx) at (1.75,-5) {$\dfrac{\partial^2 f}{\partial x\,\partial y} = 1$};
  \node[box=cgreen] (fyy) at (5.6,-5) {$\dfrac{\partial^2 f}{\partial y^2} = 2$};
  \draw[arrow, draw=cblue] (f) -- node[above left, cblue] {by $x$} (fx);
  \draw[arrow, draw=cgreen] (f) -- node[above right, cgreen] {by $y$} (fy);
  \draw[arrow, draw=cblue] (fx) -- node[left, cblue] {by $x$} (fxx);
  \draw[arrow, draw=cgreen] (fx) -- node[right, cgreen] {by $y$} (fxy);
  \draw[arrow, draw=cblue] (fy) -- node[left, cblue] {by $x$} (fyx);
  \draw[arrow, draw=cgreen] (fy) -- node[right, cgreen] {by $y$} (fyy);
  \node[corange] at (0,-6.15) {the two mixed paths agree};
  \node[font=\Large] (H) at (0,-8.1) {$H = \begin{bmatrix} {\color{cblue}6x} & {\color{corange}1} \\[3pt] {\color{corange}1} & {\color{cgreen}2} \end{bmatrix}$};
  \draw[arrow, dashed] (fxx.south) to[out=-80, in=180] (H.west);
  \draw[arrow, dashed] (fyy.south) to[out=-100, in=0] (H.east);
  \draw[arrow, dashed] (0,-6.45) -- (H.north);
\end{tikzpicture}
\end{document}
